% Hivemind Attacks — 2 October 2026. Eight main slides plus one Q&A backup.
% Build from this directory with pdfLaTeX; evidence and caveats are in \note{}.
\documentclass[aspectratio=169]{beamer}
\usetheme{Chalmers}
\usepackage[utf8]{inputenc}
\usepackage[T1]{fontenc}
\usepackage[british]{babel}
\usepackage{graphicx}
\usepackage{array}
\usepackage{tikz}
\usepackage{pgfpages}
\usetikzlibrary{positioning,arrows.meta,shapes.geometric,calc}

\title{Hivemind Attacks}
\subtitle{Two failures in aggregation coordination}
\author{Sarat Yellapragada}
\institute{Chalmers University of Technology}
\date{2 October 2026}
\shorttitle{\ChalmersSectionName}
\chalmersheadline{logo}
\chalmerssectiondividers{none}
\chalmersfooterpages{both}
\AtBeginSection[]{}
\renewcommand{\Chalmersinsertpagenumber}{\insertframenumber\,/\,\insertmainframenumber}
\setbeamerfont{frametitle}{size=\fontsize{19}{22}\selectfont,series=\bfseries}
\setlength{\Chalmersframetitlebottom}{0mm}
\ifdefined\presenternotes
  \setbeameroption{show notes on second screen=right}
\else
  \setbeameroption{hide notes}
\fi

\colorlet{attack}{ChalmersRed}
\colorlet{safe}{ChalmersGreen}
\colorlet{dark}{ChalmersBlack}
\colorlet{quiet}{ChalmersGrey}
\tikzset{
  >={Stealth[length=2.5mm]},
  every node/.style={font=\small,text=dark},
  panel/.style={rounded corners=2mm,draw=ChalmersPurple!35,fill=ChalmersPurple!6,
    align=center,inner sep=3mm},
  attackpanel/.style={panel,draw=attack!75,fill=attack!8},
  peer/.style={circle,draw=ChalmersPurple,fill=ChalmersPurple!9,
    minimum size=10mm,font=\small\bfseries},
  badpeer/.style={peer,draw=attack,fill=attack!12,text=attack},
  link/.style={->,thick,draw=ChalmersPurple},
  badlink/.style={->,thick,draw=attack},
}
\newenvironment{canvas}{\centering\begin{tikzpicture}[x=1cm,y=1cm]
  \path[use as bounding box] (0,0) rectangle (14.4,5.6);
}{\end{tikzpicture}}

\begin{document}

\begin{chalmerscover}
\begin{canvas}
  \node[anchor=west,text=white,font=\fontsize{31}{35}\selectfont\bfseries]
    at (0.25,4.25) {Hivemind Attacks};
  \node[anchor=west,text=white,font=\large] at (0.25,3.35)
    {Who joins the group? Who decides the averaging work?};
  \node[anchor=west,text=white,align=left,font=\small] at (0.25,0.95)
    {Sarat Yellapragada\\[1mm]Secura Lab, Chalmers University of Technology\\[1mm]2 October 2026};
\end{canvas}
\note{This is a roughly twelve-minute talk about two separately demonstrated control-plane attacks.
The first changes who is seen during group discovery. The second changes the ordered group
roster after an attacker is already leader. Neither experiment was run on a public WAN, and
their composition has not been measured.}
\end{chalmerscover}

\section{The coordination path}
\begin{frame}{Control messages decide the group}
\begin{canvas}
  \node[panel,minimum width=2.8cm,minimum height=1.65cm] (ads) at (1.85,3.30)
    {DHT\\advertisements};
  \node[panel,minimum width=2.8cm,minimum height=1.65cm] (view) at (5.45,3.30)
    {Peer's local view\\and leader choice};
  \node[panel,minimum width=2.8cm,minimum height=1.65cm] (roster) at (9.05,3.30)
    {Leader sends an\\ordered roster};
  \node[panel,minimum width=2.8cm,minimum height=1.65cm] (reduce) at (12.65,3.30)
    {All-reduce\\averages tensors};
  \draw[link] (ads)--(view);
  \draw[link] (view)--(roster);
  \draw[link] (roster)--(reduce);
  \node[attackpanel,minimum width=4.0cm] at (4.1,1.30)
    {Attack 1: control the lookup replies};
  \node[attackpanel,minimum width=4.0cm] at (10.3,1.30)
    {Attack 2: send conflicting rosters};
  \draw[badlink] (4.1,1.7)--(4.1,2.35);
  \draw[badlink] (10.3,1.7)--(10.3,2.35);
\end{canvas}
\note{Hivemind peers advertise availability under a group key in the DHT. A peer reads the
advertisements, makes its local leader choice, then negotiates membership. The leader sends
BEGIN\_ALLREDUCE on each join stream. The ordered roster determines which peer reduces which
shard. The map in research/control\_data\_map.pdf places the two selected attacks at checkpoints
7/13 and 14/16. This is a control-flow overview, not a claim that the attacks were combined.}
\end{frame}

\section{Attack 1: targeted matchmaking eclipse}
\begin{frame}{Attack 1: give one victim an empty view}
\begin{canvas}
  \node[panel,minimum width=3.15cm,minimum height=1.1cm] at (2.0,4.55)
    {Honest peers advertise\\under a known group key};
  \node[attackpanel,minimum width=4.15cm,minimum height=1.35cm] (dht) at (7.10,4.40)
    {Crafted DHT identities answer\\the victim's lookup};
  \node[peer] (v) at (2.10,2.15) {V};
  \node[panel,minimum width=3.25cm,minimum height=1.0cm] (empty) at (7.10,2.15)
    {``No candidates found''};
  \node[badpeer] (a) at (12.10,2.15) {A};
  \draw[link] (2.0,4.0) -- (dht.west);
  \draw[badlink] (dht.south) -- (empty.north);
  \draw[badlink] (empty.west) -- node[above,font=\footnotesize] {filtered reply} (v.east);
  \draw[badlink] (a.west) -- node[above,font=\footnotesize] {malicious peers join} (8.95,2.15);
  \node[font=\small\bfseries,text=attack,align=center] at (7.20,0.70)
    {Victim waits as leader; its group contains no other honest peer.};
\end{canvas}
\note{The DHT-only Sybils are not training participants. Their self-chosen node IDs are
placed near the known matchmaking key; the attack filters honest advertisements from the
victim's replies. An empty local list makes the victim wait as leader. Separate malicious
training peers then join. No honest peer other than the victim enters a captured group.
The attacker must control the lookup responders that have a cached copy, not just the five
record replicas. Source: research/HANDOFF.md, section 4.6; research/CONTROL\_DATA\_AUDIT.md C7.S.}
\end{frame}

\begin{frame}{Eclipse: a lookup-width threshold}
\begin{canvas}
  \node[anchor=west,font=\small\bfseries] at (0.45,5.05) {Victim's lookup width};
  \node[anchor=west,font=\small\bfseries] at (4.25,5.05) {Crafted DHT IDs};
  \node[anchor=west,font=\small\bfseries] at (8.1,5.05) {Captured victim rounds};
  \draw[ChalmersPurple!35] (0.45,4.73)--(13.95,4.73);
  \foreach \yy/\width/\ids/\result/\bar in {
    4.13/20/18/{4/30}/0.52,
    3.22/20/20/{20/20}/4.00,
    2.31/40/25/{9\%}/0.36,
    1.40/40/45/{100\%}/4.00} {
    \node[anchor=west,font=\small] at (0.45,\yy) {$k=\width$};
    \node[anchor=west,font=\small] at (4.25,\yy) {\ids};
    \fill[attack!75] (8.1,\yy-0.17) rectangle ++(\bar,0.35);
    \node[anchor=west,font=\small\bfseries] at (12.40,\yy) {\result};
  }
  \node[anchor=west,font=\footnotesize,text=quiet] at (0.45,0.47)
    {Local swarm; eight training peers, group size four, cached records.};
\end{canvas}
\note{At default bucket\_size = 20, 18 crafted identities captured 4 of 30 observed victim
rounds, while 20 captured 20 of 20. At bucket\_size = 40, 25 identities captured 9 percent
and 45 captured 100 percent in the recorded runs. The run lengths differ, so these are
observed fractions, not estimated public-network probabilities. The campaign used eight
training peers, three of them malicious, group size four, plus DHT-only identities. Five
crafted IDs hold the replicas, but cached copies from other lookup responders defeat total
filtering. Source: research/HANDOFF.md section 4.6 and raw step6\_isolation CSVs;
research/CONTROL\_DATA\_AUDIT.md C7.S.}
\end{frame}

\section{Attack 2: roster equivocation}
\begin{frame}{One group, conflicting reducer indices}
\begin{canvas}
  \node[badpeer] (leader) at (2.20,2.85) {L};
  \node[anchor=west,font=\small\bfseries,text=attack] at (0.55,1.75)
    {Malicious leader};
  \node[panel,minimum width=3.35cm] at (2.25,4.68)
    {One group ID\\Same four PeerIDs};
  \node[panel,minimum width=5.6cm] (f1) at (10.50,4.35)
    {To F1: [F1, L, F2, F3]};
  \node[panel,minimum width=5.6cm] (f2) at (10.50,2.85)
    {To F2: [F2, F1, L, F3]};
  \node[panel,minimum width=5.6cm] (f3) at (10.50,1.35)
    {To F3: [F3, F1, F2, L]};
  \draw[badlink] (leader) -- (f1.west);
  \draw[badlink] (leader) -- (f2.west);
  \draw[badlink] (leader) -- (f3.west);
  \node[anchor=west,font=\footnotesize,text=quiet] at (0.55,0.42)
    {Illustrative permutations; roster position selects the shard reducer.};
\end{canvas}
\note{This assumes the attacker is already the legitimate leader; the experiment forced
that precondition rather than measuring election frequency. The leader sends a separate
BEGIN\_ALLREDUCE to each follower. It can preserve exactly the same authenticated PeerIDs,
group ID, and gathered list length while swapping each recipient into the leader's reducer
index. Followers verify their own membership but do not cross-check a common ordered roster.
The shown permutations illustrate the mechanism; exact per-recipient payloads are in the
step17 events CSV. Source: research/CONTROL\_DATA\_AUDIT.md C14.T, lines 948 onward.}
\end{frame}

\begin{frame}{Averaging completes; outputs are wrong}
\begin{canvas}
  \node[panel,minimum width=5.70cm,minimum height=2.65cm] at (3.62,3.48) {};
  \node[font=\small\bfseries,text=ChalmersPurple] at (3.62,4.35) {Common roster};
  \node[font=\fontsize{34}{38}\selectfont\bfseries,text=safe] at (3.62,3.42) {0/5};
  \node[font=\small] at (3.62,2.65) {trials with corrupted output};
  \node[attackpanel,minimum width=5.70cm,minimum height=2.65cm] at (10.78,3.48) {};
  \node[font=\small\bfseries,text=attack] at (10.78,4.35) {Different rosters};
  \node[font=\fontsize{34}{38}\selectfont\bfseries,text=attack] at (10.78,3.42) {5/5};
  \node[font=\small] at (10.78,2.65) {trials with corrupted output};
  \node[font=\small\bfseries] at (7.20,1.25)
    {All 20 attacked peer outputs wrong. All 40 steps completed.};
  \node[font=\footnotesize,text=quiet] at (7.20,0.72)
    {Four real averagers, five paired local trials; no bans, fallbacks or timeouts.};
\end{canvas}
\note{Four real DecentralizedAverager peers used the local DHT, authenticated join streams,
and stock all-reduce. In each trial every peer began with the same tensor [0,1,2,4]; a
correct average should leave it unchanged. There were five fresh control/attack pairs.
Every control output was correct; all 20 attacked peer outputs were wrong. All 40 steps
across both arms completed, with zero reducer bans, fallback elements, RPC/all-reduce
errors, timeouts, or hard cancellations. This measures tensor corruption, not downstream
model convergence. Source: research/HANDOFF.md current thread; research/CONTROL\_DATA\_AUDIT.md
C14.T; research/results/step17\_roster\_live\_rounds\_confirm\_20260928\_confirm\_t5\_s1717.csv.}
\end{frame}

\section{Interpretation}
\begin{frame}{Two different trust boundaries}
\begin{canvas}
  \node[attackpanel,minimum width=6.2cm,minimum height=2.40cm] at (3.52,3.78) {};
  \node[font=\small\bfseries,text=attack] at (3.52,4.62) {Discovery control};
  \node[align=center,font=\small] at (3.52,3.72)
    {Crafted DHT IDs filter the victim's view\\[1mm]
     \(\rightarrow\) no other honest peer in its group};
  \node[panel,minimum width=6.2cm,minimum height=2.40cm] at (10.88,3.78) {};
  \node[font=\small\bfseries,text=ChalmersPurple] at (10.88,4.62) {Roster control};
  \node[align=center,font=\small] at (10.88,3.72)
    {An existing leader sends conflicting orders\\[1mm]
     \(\rightarrow\) silent all-reduce corruption};
  \node[panel,minimum width=12.9cm,minimum height=1.05cm] at (7.20,1.45)
    {Both were demonstrated locally. They have not been combined or tested on a public swarm.};
\end{canvas}
\note{The first attack is a selection result, not by itself a measured model-poisoning result.
The second is a silent tensor-corruption result conditional on leadership, not evidence
that the attacker can always become leader. Both use Hivemind's real local protocol path.
Neither experiment estimates success on a public WAN. A useful defence question for the
first is how to bind DHT placement and constrain a victim's lookup view. For the second,
members need agreement on one ordered roster, not merely a leader-signed message: the
leader itself is the equivocator. These defence directions have not been evaluated here.}
\end{frame}

\begin{chalmerscover}
\begin{canvas}
  \node[text=white,font=\fontsize{31}{35}\selectfont\bfseries] at (7.2,3.45)
    {Questions?};
  \node[text=white,font=\large,align=center] at (7.2,2.40)
    {Which coordination decision should be hardened first?};
\end{canvas}
\note{The two attacks answer different questions. If asked about becoming leader by
undercutting a deadline, use the backup slide and state the WAN limitations. If asked
about a combined attack, say it has not been measured. The next concrete experiment
could test roster agreement and its overhead, while a public-WAN study would need an
operator, a known group key, and permission to measure that run.}
\end{chalmerscover}

\appendix
\renewcommand{\Chalmersinsertpagenumber}{Backup}
\begin{frame}{Backup: deadline undercutting needs time}
\begin{canvas}
  \node[font=\small\bfseries] at (3.82,4.80) {5--7 s deadline window};
  \node[font=\small\bfseries] at (10.58,4.80) {11--15 s deadline window};
  \node[attackpanel,minimum width=5.45cm,minimum height=1.85cm] at (3.82,3.40)
    {\Large\bfseries 0/15 captured\\[2mm]\small with injected regional delay};
  \node[panel,minimum width=5.45cm,minimum height=1.85cm] at (10.58,3.40)
    {\Large\bfseries 15/15 captured\\[2mm]\small with injected regional delay};
  \node[panel,minimum width=12.55cm,minimum height=1.17cm] at (7.20,1.25)
    {28-node local swarm; five RIPE-derived RTT regions; known group key. Not a public-WAN result.};
\end{canvas}
\note{This is deliberately a backup, not one of the two headline attacks. Five attacker
origins times three seeds were tested per deadline window, with a 28-DHT-node localhost
swarm and injected pairwise delays drawn from RIPE Atlas ICMP anchor pings. Long-window
capture was 15/15, short-window capture 0/15; the default matchmaking time is five
seconds. The emulator omits packet loss, churn, DHT pings, connection setup, and delay
during all-reduce; it cannot establish success on a public swarm. Sources: research/meetings/
2026-10-02\_meeting\_agenda.md and research/results/step19\_wan\_delay\_rounds\_smoke\_r5\_*.csv.
The read/write timings already contain the DHT walk.}
\end{frame}

\end{document}
